\section{Restriction of concentric dyadic fans}
\label{sec:area_law_concentric_fans}

Use the maximum norm on $\mathbb R^2$. Fix a common center
$c\in\mathbb R^2$ and a natural exponent $\ell$, and put $r=2^\ell/2$.
The smaller square has origin $o_{\rm S}=c-r(1,1)$, exponent $\ell$
and cell index $(0,0)$; the larger square has origin
$o_{\rm L}=c-2r(1,1)$, exponent $\ell+1$ and the same index.
Choose identical optional midpoint subdivisions on corresponding sides,
and denote the corresponding triangles by $P_e^{\rm S}$ and $P_e^{\rm L}$.
This restriction supplies the margin in the local construction of
\cite[Section~11]{OpenAI2026AreaLaw}.

% Source: 10-geometry.tex, geometry:initial-stars, lines 333--370,
% especially 352--363, and prop:two-families, lines 299--323;
% revision adc7f1241b42e322a6451854ab7e4b4c146bf78a.

\begin{theorem}[Exact restriction of corresponding fan triangles]
    \label{thm:area_law_concentric_fan_restriction}
    \lean{TNLean.PEPS.AreaLaw.Geometry.cellFanPolygon_succ_inter_closedBall}
    \leanok
    \uses{def:area_law_cell_fan}
    For every pair of corresponding elementary bases $e$,
    \begin{align}
        P_e^{\rm L}\cap\overline B_\infty(c,r) & =P_e^{\rm S}.
    \end{align}
    The subdivisions may be arbitrary; they need only agree between
    the two fans.
\end{theorem}
\begin{proof}\leanok
    \uses{def:area_law_cell_fan,thm:area_law_cell_fan_base_radius}
    The homothety $H(y)=c+(y-c)/2$ fixes $c$ and sends the two large
    perimeter vertices to their corresponding small vertices. Affine
    maps preserve convex hulls, so $H(P_e^{\rm L})=P_e^{\rm S}$.

    A point of the larger triangle has the form
    $x=c+\lambda(y-c)$, where $0\le\lambda\le1$ and $y$ lies on its
    elementary base. The base-radius identity gives
    $\|y-c\|_\infty=2r$, hence $\|x-c\|_\infty=2r\lambda$.
    If $x\in\overline B_\infty(c,r)$, then $\lambda\le1/2$, and
    $x=c+2\lambda(H(y)-c)$ belongs to the smaller triangle.
    Conversely, halving a radial coefficient in $[0,1]$ keeps the
    resulting point in the larger triangle and within distance $r$
    of $c$. This proves both inclusions.
\end{proof}
